\documentclass[letterpaper,12pt,twoside]{book}  % book class
\usepackage[spanish,es-nodecimaldot,es-tabla]{babel} % language/hyphenation
\usepackage[utf8]{inputenc}  % input encoding
\usepackage{graphicx} % graphics library (REQUIRED by CIMATpreamble)
\usepackage[svgnames,table]{xcolor}  % color library (REQUIRED by CIMATpreamble)
\usepackage{CIMATpreamble}  % definition of the CIMAT cover page
\usepackage{mypreamble} % custom preamble
\graphicspath{{img/}}  % set path for images folder

%-----------------------------------------------------------------------------%
% Coverpage template requested by CIMAT for the master thesis in LaTeX version

%-----------------------------------------------------------------------------%
% README: Instructions & Comments
%
% > The package 'CIMATpreamble' is required for the CIMAT cover page and
%   it depends on the declared libraries above
% > To use the COVER template only fill the Cover/Titlepage block
%   with your information and make sure the '\maketitle' command is enabled
%
% > The rest of the LaTeX code was built according to the guideline provided by
%   the ShareLatex/Overleaf team:
%   "How to Write a Thesis in LaTeX (Part 1): Basic Structure"; you will find
%   the custom settings (link colors, caption specs., etc.) in the
%   'mypreamble.sty' file, so you are free to use or modify that code
% > Much of the code for the cover page is based on the Overleaf template
%   created by Salvador Pedraza Espitia:
%   https://www.overleaf.com/latex/templates/unam-dissert/mypjkyrmhgns
% > The essential idea of this code is to provide a minimal and functional
%   LaTeX document, allowing us to focus on the content of our thesis work;
%   I hope it be useful for you :) !
%
% > Credits: Kenny Yahir Mendez Ramirez, Cristal E. Mendez Ramirez
% > For comments and suggestions please send a e-mail to kenny.yahir.mz@gmail.com

%-----------------------------------------------------------------------------%
% Cover/Title page: filzzxl out with your personal information

\author{Gustavo Hernández Angeles}
\documentType{T E S I S}  % change to "T E S I N A" if it is the case
\title{Sistema de Recomendación Híbrido en Pueblos Mágicos: Integración de Análisis de Aspectos y Optimización}
\degree{Maestro en Cómputo Estadístico}
\supervisor{Dr. Norberto Alejandro Hernández Leandro}
%\supervisorSecond{Dra. M. F.}  % optional command
\cityandyear{Monterrey, Nuevo León, Xxxx de 2026}

%-----------------------------------------------------------------------------%
% Document body

\begin{document}

% Coverpage
\maketitle  % this enable the CIMAT title page!

\thispagestyle{empty}  % blank page after title page

\frontmatter

% Dedication
% \chapter*{}
% \begin{flushright}%
% 	\emph{A mi familia.} \\[2cm]

% 	\thispagestyle{empty}
% \end{flushright}

% Acknowledgements
\chapter*{Agradecimientos}
\addcontentsline{toc}{chapter}{Agradecimientos}

\noindent Expreso mi agradecimiento al Dr. Norberto Alejandro Hernández Leandro por su asesoría, así como a los doctores Alejandro Rosales Pérez y Miguel Ángel Álvarez Carmona, integrantes de mi comité, cuyo acompañamiento fue fundamental para este proyecto.\\[0.2cm]

\noindent Al Centro de Investigación en Matemáticas (CIMAT), le agradezco el haberme recibido para realizar mis estudios de posgrado y al profesorado por la sólida formación técnica recibida en diversas áreas.\\[0.2cm]

\noindent Asimismo, reconozco a la Secretaría de Ciencia, Humanidades, Tecnología e Innovación (SECIHTI) por el financiamiento otorgado, el cual permitió la conclusión exitosa de mi maestría.\\[0.2cm]

\noindent A mis padres Octavio y Abigail. Su amor y sacrificio han sido fundamentales para cumplir mis metas. Gracias por ser un ejemplo de que cualquier obstáculo puede superarse. Espero que mis logros académicos sean un reflejo de su dedicación. \\[0.2cm]

\noindent A mis hermanos, Arael y Javier. Han sido siempre un modelo a seguir para mí, sé que siempre puedo aprender más de ustedes. Gracias por siempre mostrar su apoyo y comprensión. \\[0.2cm]

\noindent A Yazmin. Tu apoyo incondicional ha sido una fuente de fortaleza para mí. Gracias por ser la mejor compañia tanto en los momentos difíciles como en los felices. \\[0.2cm]

\noindent Finalmente, a mis amigos y compañeros de clase, por hacer este camino mucho más ameno y por su apoyo constante. \\[0.2cm]

% Abstract
\chapter*{Resumen}
\addcontentsline{toc}{chapter}{Resumen}

% Keywords

% Table of contents and list of figures
\tableofcontents

\listoffigures
\listoftables

% Chapters
\mainmatter

\spacing{1}  % interline  (double) space
\setlength{\parskip}{6pt}

% put the content of your chapters in the .tex files located in the 'chapters/' folder;
% the use of '\include' instead of '\input' command could enhance the  compilation process, see the reference below

\chapter{Introducción}

El aumento exponencial de información generada por los usuarios en plataformas digitales ha transformado la manera en que las organizaciones comprenden y responden a las necesidades y preferencias de sus clientes.

\section{Trabajo previo}

\section{Hipótesis de trabajo} \label{sec:hypothesis}

\section{Objetivos} \label{sec:objective}

\section{Contribución} \label{sec:contribution}

\section{Organización de la tesis}








\chapter{Conjunto de datos} \label{chap:dataset}

En este capítulo se describe el origen, estructura y preprocesamiento del conjunto de datos utilizado en esta tesis. El corpus proviene de la tarea compartida REST-MEX 2025, organizada en el marco de IberLEF 2025. Se detalla el contexto del evento, la composición del dataset original, y el proceso de filtrado aplicado para obtener una matriz de interacciones con densidad suficiente para el entrenamiento de sistemas de recomendación.

%% ============================================================
%% EVENTO REST-MEX
%% ============================================================
\section{La tarea compartida REST-MEX} \label{sec:restmex}

REST-MEX (\textit{Researching Sentiment Evaluation in Text for Mexican Magical Towns}) es una tarea compartida organizada dentro del foro de evaluación IberLEF (\textit{Iberian Languages Evaluation Forum}), cuyo objetivo es establecer referencias comparativas (\textit{benchmarks}) para sistemas de procesamiento de lenguaje natural en el ámbito del turismo mexicano \citep{alvarez-carmona_overview_2025}.

La tarea se lanzó en 2021 con un enfoque en la predicción de polaridad de reseñas turísticas en español \citep{alvarez-carmona_overview_2021}. En ediciones subsecuentes, se incorporó la clasificación del tipo de servicio reseñado (hotel, restaurante o atractivo) \citep{alvarez-carmona_overview_2022} y la identificación del país de origen de la reseña \citep{alvarez-carmona_overview_2023}. La edición 2025 marca la cuarta iteración del evento e introduce cambios significativos: (i)~la clasificación geográfica se refina a nivel de Pueblo Mágico específico (40 pueblos predefinidos), (ii)~el corpus se amplía considerablemente, y (iii)~se fomenta el uso de técnicas avanzadas como ingeniería de \textit{prompts}, binarización de clases y selección de instancias basada en centroides.

En esta edición, 32 equipos participantes enviaron un total de 70 sistemas válidos, explorando una amplia variedad de enfoques que incluyen modelos basados en Transformers, aprendizaje por transferencia, esquemas de binarización y estrategias de selección de instancias \citep{alvarez-carmona_overview_2025}. La tarea se estructura en tres subtareas simultáneas que, dada una reseña turística, requieren determinar:

\begin{enumerate}
    \item La \textbf{polaridad} del sentimiento, en una escala ordinal de 1 (muy malo) a 5 (muy bueno).
    \item El \textbf{tipo de servicio} reseñado: \textit{Attractive}, \textit{Hotel} o \textit{Restaurant}.
    \item El \textbf{Pueblo Mágico} específico al que refiere la reseña, seleccionado de una lista predefinida de 40 destinos.
\end{enumerate}

Si bien la Tarea 4 del evento REST-MEX contempla la construcción de un sistema de recomendación para Pueblos Mágicos, en esta tesis se extiende dicho planteamiento incorporando análisis de sentimientos basado en aspectos y optimización metaheurística para la fusión de componentes, aspectos que no están contemplados en la formulación original del evento.

%% ============================================================
%% CORPUS ORIGINAL
%% ============================================================
\section{Corpus REST-MEX 2025} \label{sec:corpus}

El corpus REST-MEX 2025 consta de 297,217 escritas por turistas que visitaron destinos representativos de México y compartieron sus opiniones en la plataforma TripAdvisor entre 2002 y 2024 \citep{alvarez-carmona_overview_2025}. Cada reseña está etiquetada con una calificación de polaridad en una escala ordinal de cinco puntos $\{1, 2, 3, 4, 5\}$, el tipo de negocio (Attractive, Hotel, Restaurant) y el Pueblo Mágico correspondiente.

El corpus presenta un notable desbalance de clases en los tres ejes de clasificación. En el eje de polaridad, la clase 5 (muy bueno) concentra el $65.6\%$ de las instancias de entrenamiento ($136{,}561$ de $208{,}051$), mientras que las clases 1 y 2 representan conjuntamente menos del $5.3\%$. En el eje geográfico, tres pueblos (Tulum, Isla Mujeres y San Cristóbal de las Casas) concentran el $42.2\%$ de las instancias.

Para los fines de esta tesis, se recopiló un conjunto extendido de reseñas a partir de los archivos fuente del evento. El proceso de carga concatenó 1,559 archivos CSV, resultando en un dataset crudo de 313,836 registros con 12 columnas originales.

%% ============================================================
%% ESQUEMA Y LIMPIEZA
%% ============================================================
\section{Estructura y limpieza del dataset} \label{sec:cleaning}

Cada registro del dataset crudo contiene los siguientes campos:

\begin{table}[ht]
\centering
\caption{Esquema del dataset crudo de reseñas turísticas.}
\label{tab:schema}
\begin{tabular}{lll}
\toprule
\textbf{Campo} & \textbf{Tipo} & \textbf{Descripción} \\
\midrule
Author        & Texto     & Seudónimo del usuario en TripAdvisor \\
Titulo        & Texto     & Título de la reseña \\
Review        & Texto     & Cuerpo de la reseña \\
Calificacion  & Númerico  & Calificación numérica (1--5) \\
OrigenAutor   & Texto     & País de origen del autor \\
FechaOpinion  & Fecha     & Fecha de publicación de la reseña \\
FechaEstadia  & Fecha     & Mes y año de la estancia \\
Pueblo        & Texto     & Pueblo Mágico del establecimiento \\
Estado        & Texto     & Estado de la república \\
Pais          & Texto     & País del establecimiento \\
Tipo          & Texto     & Tipo de servicio (Hotel/Restaurant/Attractive) \\
Lugar         & Texto     & Nombre del establecimiento \\
\bottomrule
\end{tabular}
\end{table}

El proceso de limpieza consistió en:

\begin{enumerate}
    \item \textbf{Eliminación de columnas redundantes}: se descartaron \texttt{OrigenAutor} (53.5\% de valores nulos), \texttt{FechaOpinion} (redundante con \texttt{FechaEstadia}) y \texttt{Pais} (valor constante ``Mexico'').
    \item \textbf{Eliminación de registros incompletos}: se eliminaron filas con valores nulos en los campos \texttt{Author}, \texttt{Calificacion}, \texttt{Review} o \texttt{FechaEstadia}, reduciendo el dataset de $313{,}836$ a \textbf{306,854 registros}.
    \item \textbf{Conversión de fechas}: el campo \texttt{FechaEstadia}, originalmente en formato textual en español (e.g., ``enero de 2023''), se convirtió a formato \texttt{datetime} mediante un mapeo de nombres de meses en español a sus equivalentes numéricos.
\end{enumerate}

El dataset limpio resultante contiene 9 columnas y abarca estancias desde 2001 hasta 2023.

%% ============================================================
%% FILTRADO DE DENSIDAD
%% ============================================================
\section{Filtrado para densidad de interacciones} \label{sec:filtering}

Para el entrenamiento de algoritmos de filtrado colaborativo, es necesario que la matriz de interacciones usuario-ítem presente una densidad mínima que permita calcular similitudes confiables entre usuarios o entre ítems. El dataset limpio de $306{,}854$ reseñas contiene un gran número de usuarios que han reseñado establecimientos en un solo Pueblo Mágico, lo cual resulta insuficiente para la tarea de recomendación a nivel de pueblos.

Se diseñó un procedimiento de filtrado iterativo con dos criterios de densidad:

\begin{enumerate}
    \item \textbf{Usuarios activos}: retener únicamente usuarios que hayan reseñado establecimientos en al menos $\tau_u = 3$ Pueblos Mágicos distintos.
    \item \textbf{Pueblos con cobertura}: retener únicamente Pueblos Mágicos que contengan al menos $\tau_p = 10$ establecimientos (\textit{Lugares}) distintos con reseñas.
\end{enumerate}

Dado que la eliminación de usuarios puede dejar pueblos sin cobertura suficiente (y viceversa), el filtrado se aplica de forma iterativa hasta convergencia, es decir, hasta que ningún registro adicional sea eliminado. El Algoritmo~\ref{alg:filter} formaliza este procedimiento.

\begin{algorithm}[ht]
\caption{Filtrado iterativo para densidad de interacciones.}
\label{alg:filter}
\begin{algorithmic}[1]
\REQUIRE Dataset $\mathcal{D}$, umbrales $\tau_u$, $\tau_p$
\ENSURE Dataset filtrado $\mathcal{D}'$
\STATE $\mathcal{D}' \leftarrow \mathcal{D}$
\REPEAT
    \STATE $n_{\text{prev}} \leftarrow |\mathcal{D}'|$
    \STATE $\mathcal{U}_{\text{dense}} \leftarrow \{u \in \mathcal{U} \mid |\{p : \exists\, r \in \mathcal{D}' \text{ con } r.\text{Author} = u,\; r.\text{Pueblo} = p\}| \geq \tau_u\}$
    \STATE $\mathcal{D}' \leftarrow \{r \in \mathcal{D}' \mid r.\text{Author} \in \mathcal{U}_{\text{dense}}\}$
    \STATE $\mathcal{P}_{\text{dense}} \leftarrow \{p \in \mathcal{P} \mid |\{l : \exists\, r \in \mathcal{D}' \text{ con } r.\text{Pueblo} = p,\; r.\text{Lugar} = l\}| \geq \tau_p\}$
    \STATE $\mathcal{D}' \leftarrow \{r \in \mathcal{D}' \mid r.\text{Pueblo} \in \mathcal{P}_{\text{dense}}\}$
\UNTIL{$|\mathcal{D}'| = n_{\text{prev}}$}
\RETURN $\mathcal{D}'$
\end{algorithmic}
\end{algorithm}

%% ============================================================
%% DATASET RESULTANTE
%% ============================================================
\section{Dataset de trabajo} \label{sec:working_dataset}

La aplicación del filtrado iterativo con $\tau_u = 3$ y $\tau_p = 10$ reduce el dataset de $306{,}854$ a \textbf{79,264 registros}, correspondientes a \textbf{9,582 usuarios} y \textbf{48 Pueblos Mágicos}. La Tabla~\ref{tab:dataset_stats} resume las estadísticas principales antes y después del filtrado.

\begin{table}[ht]
\centering
\caption{Estadísticas del dataset antes y después del filtrado de densidad.}
\label{tab:dataset_stats}
\begin{tabular}{lrr}
\toprule
\textbf{Métrica} & \textbf{Dataset limpio} & \textbf{Dataset filtrado} \\
\midrule
Reseñas totales    & 306,854 & 79,264 \\
Usuarios únicos    & ---     & 9,582  \\
Pueblos Mágicos    & ---     & 48     \\
Mín. pueblos/usuario & 1     & 3      \\
Máx. pueblos/usuario & ---   & 28     \\
Mín. lugares/pueblo  & ---   & 10     \\
\bottomrule
\end{tabular}
\end{table}

La distribución de pueblos en el dataset filtrado presenta un marcado desbalance geográfico. Tulum es el pueblo con mayor representación ($10{,}051$ títulos únicos), seguido de Isla Mujeres ($6{,}082$) y San Cristóbal de las Casas ($4{,}902$), mientras que pueblos como Tecate ($175$), San Sebastián del Oeste ($193$) y Santiago ($242$) tienen una representación considerablemente menor.

En términos de la matriz de utilidad para el sistema de recomendación---donde las filas representan usuarios y las columnas Pueblos Mágicos---la dispersión (\textit{sparsity}) se calcula como:
\begin{equation} \label{eq:sparsity}
    \text{sparsity} = 1 - \frac{|\{(u, p) : \exists\, r \in \mathcal{D}'\}|}{|\mathcal{U}| \times |\mathcal{P}|} = 1 - \frac{n_{\text{interacciones}}}{9{,}582 \times 48} \approx 90.32\%
\end{equation}

Esta dispersión del $90.32\%$ implica que cada usuario ha interactuado, en promedio, con menos de 5 de los 48 pueblos disponibles. En cuanto a las calificaciones, la media global es $\bar{r} = 4.41$ con el $61.3\%$ de las calificaciones siendo de 5 estrellas, reflejando el sesgo positivo característico de las plataformas de reseñas turísticas. Estos desafíos---alta dispersión y desbalance en calificaciones---motivan la adopción de un enfoque híbrido que combine señales colaborativas con información semántica extraída de las reseñas textuales, como se detalla en los capítulos subsecuentes.


\chapter{Marco teórico} \label{chap:background}

En este capítulo se presentan los fundamentos teóricos que sustentan el desarrollo de esta tesis. Se abordan tres ejes principales: los sistemas de recomendación, el procesamiento de lenguaje natural con enfoque en el análisis de sentimientos basado en aspectos, y las técnicas de optimización metaheurística. Adicionalmente, se discuten las métricas de evaluación empleadas y los fundamentos de explicabilidad en sistemas de recomendación.

%% ============================================================
%% 2.1  SISTEMAS DE RECOMENDACIÓN
%% ============================================================
\section{Sistemas de recomendación} \label{sec:recsys}

% Definición formal del problema de recomendación.
% Matriz de utilidad R ∈ ℝ^{|U|×|I|}, donde la mayoría de entradas son desconocidas.
% Contextualizar con el dataset REST-MEX: |U|=9,582 usuarios, |I|=48 pueblos, sparsity=90.32%.

% --------------------------------------------------------------
\subsection{Filtrado colaborativo} \label{subsec:cf}

% Idea central: usuarios con comportamiento similar en el pasado tendrán preferencias similares en el futuro.

\subsubsection{Basado en Usuarios} \label{subsubsec:cf_memory}
% User-based CF: similaridad entre usuarios (coseno, Pearson).
% Item-based CF: similaridad entre ítems.
% Predicción como promedio ponderado por similaridad.
% Limitaciones: escalabilidad, sparsity, cold-start.

\subsubsection{Basado en Items} \label{subsubsec:cf_model}
% Factorización matricial: R ≈ P·Qᵀ, donde P ∈ ℝ^{|U|×k}, Q ∈ ℝ^{|I|×k}.
% SVD, SVD++, NMF.
% Función de pérdida con regularización:
%   min_{p,q} Σ_{(u,i)∈κ} (r_{ui} - p_u^T q_i)^2 + λ(||p_u||^2 + ||q_i||^2)
% Modelos baseline de la librería Surprise: KNNBasic, KNNWithMeans, SVD, SVD++, NMF, BaselineOnly.
% Ejemplo: en REST-MEX, con k factores latentes representando dimensiones como "cultura" o "gastronomía".

% --------------------------------------------------------------
\subsection{Filtrado basado en contenido} \label{subsec:cbf}

% Representación de ítems mediante vectores de características (perfil del ítem).
% Perfil del usuario como agregación de los perfiles de ítems consumidos.
% Similitud coseno entre perfil del usuario y candidatos.
% En este proyecto: características extraídas de reseñas textuales via ABSA (Sección \ref{sec:absa}).
% Ejemplo: perfil de aspectos de un pueblo = distribución de sentimientos por aspecto (servicio, gastronomía, naturaleza, ...).
% Ventaja sobre CF puro: mitiga cold-start de ítems.

% --------------------------------------------------------------
\subsection{Sistemas híbridos} \label{subsec:hybrid}

% Taxonomía de Burke (2002): weighted, switching, mixed, feature combination, cascade, feature augmentation, meta-level.
% Justificar la estrategia elegida para esta tesis (feature augmentation / weighted).
% Formalización: r̂_{ui} = α · r̂_{ui}^{CF} + (1-α) · r̂_{ui}^{CBF}, donde α es el parámetro de fusión.
% Extensión: α puede ser un vector de pesos optimizado por metaheurística (Sección \ref{sec:metaheuristic}).
% Ventaja: combina señales complementarias (patrones de rating + semántica textual).

% --------------------------------------------------------------
\subsection{Problema del arranque en frío y dispersión} \label{subsec:cold_start}

% Cold-start de usuario: nuevos usuarios sin historial → CBF puede explotar contenido textual.
% Cold-start de ítem: nuevos pueblos sin ratings → ABSA genera perfil desde reseñas.
% Sparsity: 90.32% en REST-MEX. Impacto en CF basado en memoria.
% Estrategias de mitigación: factorización matricial, incorporación de side information, sistemas híbridos.


%% ============================================================
%% 2.2  PROCESAMIENTO DE LENGUAJE NATURAL
%% ============================================================
\section{Procesamiento de lenguaje natural} \label{sec:nlp}

% Breve contexto: el texto de reseñas como fuente de información semántica complementaria a los ratings numéricos.

% --------------------------------------------------------------
\subsection{Representación de texto} \label{subsec:text_repr}

% Bag-of-Words, TF-IDF → limitaciones (pérdida de orden, alta dimensionalidad).
% Word embeddings (Word2Vec, GloVe) → representaciones densas.
% Embeddings contextuales (BERT, RoBERTa) → representaciones dependientes del contexto.

% --------------------------------------------------------------
\subsection{Arquitectura Transformer} \label{subsec:transformers}

% Mecanismo de self-attention: Attention(Q,K,V) = softmax(QKᵀ/√d_k)V.
% Multi-head attention, codificación posicional.
% Modelos pre-entrenados y fine-tuning para tareas downstream.
% Transfer learning: ventaja en dominios con datos etiquetados escasos (como turismo en español).


%% ============================================================
%% 2.3  ANÁLISIS DE SENTIMIENTOS BASADO EN ASPECTOS (ABSA)
%% ============================================================
\section{Análisis de sentimientos basado en aspectos} \label{sec:absa}

% Definición formal: dado un texto t, identificar un conjunto de tuplas {(a_i, s_i)} donde a_i es un aspecto y s_i ∈ {positivo, negativo, neutro}.
% Diferencia con análisis de sentimientos general: granularidad a nivel de aspecto, no de documento.
% Ejemplo REST-MEX: "La comida estaba deliciosa pero el hotel tenía mal servicio"
%   → {(gastronomía, positivo), (servicio, negativo)}.

% --------------------------------------------------------------
\subsection{Subtareas del ABSA} \label{subsec:absa_subtasks}

% Aspect Term Extraction (ATE): identificar menciones explícitas de aspectos en el texto.
% Aspect Category Detection (ACD): asignar categorías predefinidas (e.g., servicio, gastronomía, naturaleza).
% Aspect Sentiment Classification (ASC): determinar polaridad por aspecto.
% En esta tesis: ACD vía Zero-Shot + ASC vía modelo de sentimiento fine-tuned.

% --------------------------------------------------------------
\subsection{Clasificación Zero-Shot} \label{subsec:zero_shot}

% Paradigma: clasificar sin ejemplos etiquetados para las clases objetivo.
% Basado en Natural Language Inference (NLI): premisa=texto, hipótesis="Este texto habla sobre {aspecto}".
% Modelo utilizado: Recognai/zeroshot_selectra_medium (SELECTRA, entrenado en español).
% Ventaja: no requiere datos etiquetados específicos del dominio turístico para aspectos.
% Limitación: dependencia de la formulación de las hipótesis y del conjunto de labels candidatos.

% --------------------------------------------------------------
\subsection{Análisis de sentimientos con modelos pre-entrenados} \label{subsec:sentiment_models}

% Modelos para español: BETO, RoBERTuito, MarIA.
% Modelo utilizado: pysentimiento/robertuito-sentiment-analysis.
%   - Pre-entrenado en ~500M tweets en español.
%   - Fine-tuned para clasificación de sentimiento (POS, NEG, NEU).
% Pipeline: texto → segmentación en chunks → clasificación de aspecto (Zero-Shot) → clasificación de sentimiento → agregación.


%% ============================================================
%% 2.4  OPTIMIZACIÓN METAHEURÍSTICA
%% ============================================================
\section{Optimización metaheurística} \label{sec:metaheuristic}

% Problema: encontrar los pesos óptimos de fusión del sistema híbrido.
% Formulación: max_{θ} f(θ), donde θ = (α, hiperparámetros del modelo) y f es una métrica de ranking (e.g., NDCG@K).
% Espacio de búsqueda no convexo, no diferenciable → justificación de metaheurísticas.

% --------------------------------------------------------------
\subsection{Optimización por enjambre de partículas (PSO)} \label{subsec:pso}

% Algoritmo inspirado en comportamiento colectivo (bandadas, cardúmenes).
% Partícula i con posición x_i y velocidad v_i:
%   v_i(t+1) = ω·v_i(t) + c₁·r₁·(pbest_i - x_i(t)) + c₂·r₂·(gbest - x_i(t))
%   x_i(t+1) = x_i(t) + v_i(t+1)
% Parámetros: inercia ω, coeficientes cognitivo c₁ y social c₂.
% Balance exploración-explotación.

% --------------------------------------------------------------
\subsection{Optimización bayesiana} \label{subsec:bayesian_opt}

% Modelo sustituto (surrogate): Proceso Gaussiano GP(μ, k) que aproxima f(θ).
% Función de adquisición: Expected Improvement, UCB, Probability of Improvement.
%   EI(θ) = E[max(f(θ) - f(θ⁺), 0)]
% Ventaja: eficiente en evaluaciones costosas (cada evaluación requiere entrenar y evaluar el RecSys completo).
% Ciclo: evaluar → actualizar GP → optimizar adquisición → siguiente punto.

% --------------------------------------------------------------
\subsection{Evolución diferencial} \label{subsec:de}

% Operadores: mutación, cruce, selección.
% Mutación: v_i = x_{r1} + F·(x_{r2} - x_{r3}), donde F ∈ [0,2] es el factor de escala.
% Cruce binomial con probabilidad CR.
% Selección greedy: x_i(t+1) = u_i si f(u_i) ≥ f(x_i(t)), else x_i(t).
% Variantes: DE/rand/1, DE/best/1, DE/current-to-best/1.


%% ============================================================
%% 2.5  MÉTRICAS DE EVALUACIÓN
%% ============================================================
\section{Métricas de evaluación} \label{sec:metrics}

% Contexto: evaluación de sistemas Top-K; no basta RMSE/MAE del rating prediction.

% --------------------------------------------------------------
\subsection{Métricas de ranking} \label{subsec:ranking_metrics}

\subsubsection{Precision@K y Recall@K}
% Precision@K = |{ítems relevantes en top-K}| / K
% Recall@K = |{ítems relevantes en top-K}| / |{total ítems relevantes}|
% Trade-off precision-recall según K.

\subsubsection{NDCG@K}
% DCG@K = Σ_{i=1}^{K} (2^{rel_i} - 1) / log₂(i+1)
% NDCG@K = DCG@K / IDCG@K
% Penaliza ítems relevantes en posiciones bajas del ranking.
% Métrica principal en esta tesis.

\subsubsection{MRR (Mean Reciprocal Rank)}
% MRR = (1/|U|) Σ_{u=1}^{|U|} 1/rank_u
% Evalúa la posición del primer ítem relevante.

% --------------------------------------------------------------
\subsection{Métricas beyond-accuracy} \label{subsec:beyond_accuracy}

\subsubsection{Diversidad}
% Intra-list diversity: ILD = (2 / K(K-1)) Σ_{i<j} dist(item_i, item_j)
% Importancia en turismo: evitar recomendar K pueblos del mismo estado/tipo.

\subsubsection{Cobertura}
% Catalog coverage = |{ítems recomendados al menos una vez}| / |I|
% Relevante dado el desbalance geográfico del dataset (algunos pueblos con muchas más reseñas).


%% ============================================================
%% 2.6  EXPLICABILIDAD EN SISTEMAS DE RECOMENDACIÓN
%% ============================================================
\section{Explicabilidad en sistemas de recomendación} \label{sec:xai}

% Motivación: confianza del usuario, transparencia, regulación.
% Tipos de explicación: feature-based, example-based, textual.
% En esta tesis: explicabilidad basada en aspectos.
%   Ejemplo: "Se te recomienda Taxco porque valoras positivamente la cultura y la gastronomía,
%             aspectos altamente evaluados en este pueblo."
% Conexión ABSA → explicabilidad: los aspectos extraídos sirven como justificación interpretable.
% Post-hoc vs inherent interpretability.


\chapter{Metodología y Resultados}\label{chap:eda}


\include{chapters/chap_04}

\include{chapters/conclusions}

% Bibliography
\bibliographystyle{apacite}  % APA style according to apacite package: see mypreamble.sty line 21
\bibliography{references}

% Appendix
\appendix
\include{chapters/appendix}

\backmatter

\end{document}

% LaTeX references (technical issues, workarounds)

% https://www.overleaf.com/learn/latex/Management_in_a_large_project ('including files' section)
% https://tex.stackexchange.com/questions/64839/how-to-change-listing-caption
% https://stackoverflow.com/questions/2709898/change-list-of-listings-text
% https://tex.stackexchange.com/questions/664/why-should-i-use-usepackaget1fontenc
